\subsection{Theta jets and the lower bound in the thin region}
\label{shared-actual-thin-bergman}

The issue in the thin region is quantitative nonvanishing of a
prescribed determinant differential operator.  A short spanning set
does not give this.  We construct a section of the original period
torus, and then use the finite theta basis to obtain an actual cusp
section~\cite{bergman-thin-source}.  Maximizing a determinant over a bounded set of lattice
frequencies avoids any assumption on the condition number of their
projections to one complex embedding.

\paragraph{The period torus and its parameters.}
Let $E$ be a CM field of degree $2f$, put $R=\operatorname{rd}(E)$,
and let
\[
 V=\bigoplus_{\tau:E\hookrightarrow\mathbb C}\mathbb C^{s_\tau},
 \qquad s_\tau+s_{\bar\tau}=r,\qquad g=\dim_{\mathbb C}V=fr.
\]
Use the standard Hermitian form $H$, linear in its first argument.
The original period lattice $\Gamma\subset V$ is stable under
$\mathcal O_E$, and $\omega=\operatorname{Im}H$ is integral and
unimodular on $\Gamma$.  Its semicharacter
$\varepsilon:\Gamma\to\{1,-1\}$ satisfies
\[
 \varepsilon(\gamma+\eta)
   =\varepsilon(\gamma)\varepsilon(\eta)e^{\pi i\omega(\gamma,\eta)}.
\]
We retain its original isotropic normalization:
$\varepsilon(\gamma)=1$ whenever
$\omega(a\gamma,\gamma)=0$ for every $a\in\mathcal O_E$.
Write $\lambda=\min_{0\ne\gamma\in\Gamma}\|\gamma\|$.

For an integer $q\ge1$, let $\mathcal H_q$ consist of the entire
functions with transformation law
\[
 F(z+\gamma)=\varepsilon(\gamma)^q
  e^{\pi q H(z,\gamma)+\pi q\|\gamma\|^2/2}F(z).
\]
Its squared norm is the integral of
$|F(z)|^2e^{-\pi q\|z\|^2}$ on $V/\Gamma$.
This is also the probability norm: unimodularity of $\omega$
gives Euclidean covolume one.  On $V^n$ use the exterior tensor
power of this line and the product probability measure.

Let $P$ be a product of powers of determinants.  Each determinant
uses $i\le s_\tau$ rows at one embedding $\tau$, and $i$
distinct columns among the $n$ copies of $V$.  Its exponent is an
integer $c\ge1$.  Different determinants have disjoint column
sets.  Write $D=\sum ic\ge1$ for its degree.  The rows can be any
specified orthonormal coordinate rows.  The differential operator
$P(\partial)$ uses ordinary derivatives, with no factorial
inserted into its definition.

Choose $m\ge1$ with $s_\tau\le m$ at at least one embedding, and
put
\begin{gather*}
 d_*=2fm,\qquad \Omega=2f^{3/2}R,\qquad
 \theta^2=\frac{fR_0^2R^{21}}q,\qquad
 b_0=\frac q{R_0^2R^{28}},\qquad t_0=qb_0,\\
 C=8100m^4d_*^4R^7,\qquad
 \phi=\begin{cases}2f/(R^2\lambda^2),&\lambda<\theta,\\
                    0,&\lambda\ge\theta.
       \end{cases}
\end{gather*}
The following follow from the original numerical hypotheses and rank
conditions:
\begin{gather}
 r\ge2,\quad m^4\le f,\quad f^2\le R,\quad R\ge16,
 \quad q\ge6,\quad R_0^2\ge90(g+D),\label{eq:thin-jet-parameters-a}\\
 \Omega\theta^2<1,\qquad b_0\ge1,\qquad
 t_0\ge4D\bigl(7+\log(qD)\bigr)+d_*+1.
 \label{eq:thin-jet-parameters-b}
\end{gather}
Indeed the original ranks have $r=k+m-1$, $m\ge2$ and
$k\ge m+3$, so $r\ge6$.  The original inequalities include
$d_*^2\Omega\theta^2<1$ and $90m^2d_*^2R\le q/2$.
Since $d_*=2fm\ge2$, they imply $\Omega\theta^2<1$ and
$q\ge180m^2d_*^2R\ge11520$.  The other displayed inequalities
are retained verbatim.  In particular this argument does not infer
$r\ge2$ merely from $fr\ge2$.

\begin{theorem}[The required determinant jet]
\label{thm:actual-thin-determinant-jet}
Under these hypotheses, there is a section $F$ of the original
line on $(V/\Gamma)^n$ such that
\begin{equation}\label{eq:thin-jet-lower}
 \frac{|P(\partial)F(0)|^2}{\|F\|_2^2}
   \ge 2^{-3n}16^{-D}
                  \exp\left(-\frac{2\pi nC\phi}{q}\right).
\end{equation}
In particular this implies the prescribed lower bound
\[
 (4K_6)^{-n}2^{-(r+2)D}
                  \exp\left(-\frac{2\pi nC\phi}{q}\right),
 \qquad K_6=6(2d_*)^{d_*/2}.
\]
The theorem concerns the original principally polarized period
lattice, including its semicharacter and CM action.
\end{theorem}

\begin{proof}
\emph{An isotropic sublattice and its dual frequencies.}
There are $2f$ rationally independent elements of
$\mathcal O_E$ whose absolute values at all embeddings are at
most $\Omega$.  Indeed the trace-dual minimum is at least
$\sqrt{2f}/R$; the transference argument proved in
Section~\ref{sec:actual-short-theta-generators} gives Euclidean
length at most $(2f)^{3/2}R$.  A single CM embedding accounts for
two equal summands in its squared Euclidean norm.

Start with $U=0$ and $I=\Gamma\cap U$.  We build an
$E$-stable, real isotropic subspace $U$ with the following property:
if $b$ is the vector in $U$ representing the functional
$v\mapsto\omega(\gamma,v)$ on $U$, for $\gamma\in\Gamma$,
then
\begin{equation}\label{eq:thin-dual-separation}
             b\ne0\quad\Longrightarrow\quad
                         \|b\|\ge\frac1{\Omega\theta}.
\end{equation}
Suppose $x\in\Gamma\setminus U$ has distance less than
$\theta$ from $U$.  The property forces $\omega(x,U)=0$,
since $\Omega\theta^2<1$.  Write $x=u+w$, with $u\in U$;
then $w$ is orthogonal to $U+iU$ and $\|w\|<\theta$.
For each of the short rational basis elements $a$, the integer
$\omega(ax,x)=\omega(aw,w)$ has absolute value less than one.
It is zero.  Rational linearity gives the same conclusion for
every $a\in E$.  Thus adjoining the $E$-span of $x$ preserves
isotropy.

Property~\eqref{eq:thin-dual-separation} is preserved as well.
If a pairing is already nonzero on the old $U$, use its old
orthogonal projection.  Otherwise some pairing with $ax$, for
one of the short rational basis elements $a$, is a nonzero
integer.  The old component does not contribute, so its absolute
value is at most $\|b\|\|aw\|\le\|b\|\Omega\theta$.
This proves the asserted separation.

The process stops after at most $m$ steps.  More explicitly, the
complexification of an isotropic $E$-subspace of rank $\rho$
injects into $V$: $U\cap iU=0$.  At every embedding its
complex dimension is $\rho$, whence $\rho\le s_\tau$ for
every $\tau$.  At termination put $d=\dim_{\mathbb R}U=2f\rho
\le d_*$.  The lattice $I=\Gamma\cap U$ is saturated, and
\begin{equation}\label{eq:thin-quotient-separation}
 x\in\Gamma\setminus I\quad\Longrightarrow\quad
                         \operatorname{dist}(x,U)\ge\theta.
\end{equation}
The projected vectors used at each step, multiplied by the $2f$
short coefficient elements, also give a full sublattice of $I$
with successive Gram--Schmidt lengths at most $\Omega\theta$.
Consequently
\begin{equation}\label{eq:thin-isotropic-covolume}
                         \operatorname{covol}(I)
                                      \le(\Omega\theta)^d.
\end{equation}
The semicharacter is one on $I$ by its retained normalization.
The space $U$ is zero exactly when $\lambda\ge\theta$.

Write $W=U+iU$ and $A=W^\perp$.  Each period is written
\[
                   \gamma=t+ib+a,
             \qquad t,b\in U,\quad a\in A.
\]
Its frequencies $b$ constitute exactly the Euclidean dual lattice
$L=I^*$ in $U$.  A homomorphism $I\to\mathbb Z$ extends to
$\Gamma$ because $I$ is saturated, and unimodularity of
$\omega$ then supplies its representing period.  Thus this is
an equality of lattices, not an inclusion of unspecified index.
The periods with $b=0$, modulo $I$, project to a lattice in $A$
with minimum at least $\theta$.  Every other frequency fibre is
a translate of that lattice.  By
\eqref{eq:thin-dual-separation},
\begin{equation}\label{eq:thin-frequency-minimum}
 \lambda_L\ge(\Omega\theta)^{-1}\ge\sqrt{b_0}.
\end{equation}
The second inequality is equivalent to $f\Omega^2\le R^7$,
which follows from $f^2\le R$ and $R\ge16$.

\emph{A bounded set controlling every relevant determinant.}
Let $\mathcal V$ be the closed Voronoi cell of $L$ and put
$S=L\cap3\mathcal V$.  Then
\begin{equation}\label{eq:thin-voronoi-convex}
                            \mathcal V\subset\operatorname{conv}S.
\end{equation}
To see this, for $v\in\mathcal V$ choose $\ell\in L$ with
$2v-\ell\in\mathcal V$.  Then $\ell\in3\mathcal V$ and
$v=(\ell+(2v-\ell))/2$.  Iterate this identity.  Its remainders
stay in the compact cell and their coefficients tend to zero;
the finite set $S$ includes zero, so the limit belongs to its
closed convex hull.

For $v\in U$ its Voronoi gauge is
\[
 \mu(v)=\sup_{0\ne\ell\in L}
                  \frac{2|\langle v,\ell\rangle|}{\|\ell\|^2}
                      \le\frac{2\|v\|}{\lambda_L}.
\]
Thus $v$ belongs to $\mu(v)\operatorname{conv}S$.
Moreover the covering radius of $L$ is at most
$d\lambda_d(L)/2\le d^3/(2\lambda)$, by transference and
$\lambda_1(I)\ge\lambda$.  Every $u\in S$ consequently has
\begin{equation}\label{eq:thin-frequency-radius}
 \|u\|^2\le\frac{9d^6}{4\lambda^2}\le C\phi.
\end{equation}
For the last inequality substitute $d\le2fm$ and
$\phi=2f/(R^2\lambda^2)$; it is implied by
$f\le1800m^2R^5$.  When $U=0$, set $S=\{0\}$ and only
use the abelian directions below.

Fix a block of size $i$ and exponent $c$ at $\tau$.  Consider
the projections to $V_\tau$ of $S$, together with the ball of
radius $\sqrt{qc}$ in $A\cap V_\tau$.  Maximize the absolute
value of the specified $i$-row determinant on $i$ choices from
this compact set, and call its maximum $\Delta$.  The maximum
controls determinants with columns in the complex balanced convex
hull as well, by multilinearity and the triangle inequality.
Give this span the norm whose unit ball is that hull.  For symmetric
tensors we use its projective tensor norm: the infimum of
$\sum_j\prod_k\|v_{jk}\|$ over representations as sums of
symmetrized pure tensors.  Symmetrization is the average over
permutations.  This specifies the tensor estimates used below.

The real form $U$ pairs $W\cap V_\tau$ with
$W\cap V_{\bar\tau}$ by conjugation.  A unit vector in the
first summand is the projection of a vector in $U$ of norm
$\sqrt2$.  Since $\mathcal V$ contains the radius
$\lambda_L/2$ ball, \eqref{eq:thin-voronoi-convex} and
$\lambda_L\ge1$ show that this balanced convex hull contains
the radius $1/4$ ball in all of $V_\tau$.  Hence
\begin{equation}\label{eq:thin-leading-determinant}
                              \Delta\ge4^{-i}.
\end{equation}
No lower singular-value bound for a chosen lattice basis is used.

\emph{Actual Appell--Humbert sections.}
The real form $U\subset W$ defines the complex bilinear form
$B$ whose restriction to $U$ is its Euclidean form.  Conjugation
in this real form interchanges $\tau$ and $\bar\tau$.
Consequently $B(z,z)=0$ on every one-embedding subspace
$V_\tau$.  For $u\in L$ set
\[
 e_u(z)=\exp\left(\frac{\pi q}2B(z_W,z_W)
                                  +2\pi iB(z_W,u)\right).
\]
It has the prescribed transformation law on $I$.
For a toric column selected above use
$h=e_u/(2\pi i)^c$.  For an abelian column
$\sqrt{qc}\,\xi$, where $\xi\in A\cap V_\tau$ is a unit
vector, use
\[
                       h(z)=e_0(z)
                     \frac{(\sqrt{qc}\,H(z,\xi))^c}{c!}.
\]
For an unused column use $h=e_0$.
Periodize these functions by the original lattice:
\begin{equation}\label{eq:thin-periodization}
 \Theta_h(z)=\sum_{\gamma\in\Gamma/I}
   \varepsilon(\gamma)^q
   e^{\pi qH(z,\gamma)-\pi q\|\gamma\|^2/2}h(z-\gamma).
\end{equation}
The summand is unchanged by changing the representative modulo
$I$.  The semicharacter identity makes the translation operators
a representation of $\Gamma$, so the sum belongs to
$\mathcal H_q$.  To justify the sum, choose the real $U$
coordinate of each representative in a bounded fundamental
parallelepiped.  In the remaining coordinates its modulus and
every derivative on a compact set are bounded by a polynomial
times a decaying Gaussian.  The frequency lattice and the
separated affine lattices in $A$ give locally uniform absolute
convergence.  This also justifies differentiation term by term.

For $\gamma=t+ib+a$, the scalar in the toric summand at zero
has modulus
\begin{equation}\label{eq:thin-summand-modulus}
  \exp\left(-\pi q\|b\|^2-\frac{\pi q}2\|a\|^2
                                          +2\pi\langle b,u\rangle\right).
\end{equation}
Because $u\in3\mathcal V$, this is at most
$\exp(-\pi(q-3)\|b\|^2-\pi q\|a\|^2/2)$.
Along $V_\tau$ its linear exponential is exactly
\[
            2\pi iH(z,u-qb)+\pi qH(z,a).
\]
The toric $c$-jet, after its normalization, is therefore the
$c$-th power of the functional with column
$u-qb-(q/(2i))a$, with the conjugation in the second argument
of $H$ understood.  The identity summand gives the selected column
$u$.

For the abelian column put $v=\sqrt{qc}\,\xi$.  Its $c$-jet
is the sum, for $0\le l\le c$, of the symmetric tensors
\begin{equation}\label{eq:thin-abelian-jet}
       \frac{\binom cl}{l!}(-H(a,v))^l
                  H(\,\cdot\,,v)^{c-l}\,\mathcal L_\gamma^l,
 \qquad \mathcal L_\gamma(z)=-2\pi iqH(z,b)+\pi qH(z,a),
\end{equation}
times its scalar Gaussian.  This identity follows by expanding
the polynomial and exponential before taking the $c$-th
ordinary derivative.  At $\gamma=0$ only the chosen column
$v$ remains.  The factorial in this formula is essential.

\emph{Uniform bounds for the error terms.}
Here are the estimates including their constants.  Put
\[
              A_0=q\theta^2\ge90(g+D),\qquad B_0=qb_0=t_0.
\]
For a nonzero frequency put $x=\|b\|/\lambda_L\ge1$, and
put $y=\sqrt q\|a\|$.  Measure a column by the balanced convex
hull used to define $\Delta$.  A toric jet has total tensor
coefficient norm at most
$(1+2qx+y/(2\sqrt c))^c$.
Formula~\eqref{eq:thin-abelian-jet} bounds the abelian coefficient
norm by
\[
 \sum_{l=0}^c\frac{\binom cl}{l!}
                  (4\pi q\sqrt c\,xy+\pi y^2)^l.
\]
Both are at most
\begin{equation}\label{eq:thin-nonzero-frequency-polynomial}
                    (20q\sqrt c\,x)^c(1+y)^{2c}
                         \quad(b\ne0).
\end{equation}
Indeed $1+4\pi q\sqrt c\,xy+\pi y^2$ is at most
$4(1+4q\sqrt c\,x)(1+y)^2$, and $l!\ge1$.
For $b=0$ the sharper bound is
\begin{equation}\label{eq:thin-zero-frequency-polynomial}
                         \exp(2\sqrt{\pi c}\,y).
\end{equation}
For the abelian jet use $\binom cl\le c^l/l!$ and
$\sum_{l\ge0}t^l/(l!)^2\le e^{2\sqrt t}$; the latter follows
by comparing with the even terms of the exponential using
$\binom{2l}l\le4^l$.  For the toric jet use
$\log(1+y/(2\sqrt c))\le y/(2\sqrt c)$.

We record the elementary packing bounds used here.  In a real
lattice of dimension $a$ and minimum $\delta$, at most
$(1+2R/\delta)^a$ points lie in a ball of radius $R$; this
also holds for any translate.  Split into shells of width
$\delta/2$.  Since $A_0\ge90(g+D)$, for every affine lattice
in $A$ with minimum at least $\sqrt{A_0}$ one obtains
\[
                         \sum_z e^{-\pi\|z\|^2/4}\le2.
\]
There is at most one point in the open radius $\sqrt{A_0}/2$
ball.  The remaining shells are bounded by
$\sum_{j\ge1}(j+2)^{2g}e^{-\pi A_0j^2/16}<1$.
For the lattice itself, excluding zero, the same shell calculation
starts at $j=2$ and gives
\[
                  \sum_{z\ne0}e^{-\pi\|z\|^2/4}
                         \le4\,4^{2g}e^{-\pi A_0/4}.
\]
For example $j+2\le4e^{(j^2-4)/8}$ for $j\ge2$,
and $j+2\le3e^{(j^2-1)/6}$ for $j\ge1$ prove these displayed
bounds by a geometric series.

When $b=0$ and $a\ne0$, $y\ge\sqrt{A_0}$, and
$2\sqrt{\pi c}\,y\le\pi y^2/4$.
Thus \eqref{eq:thin-zero-frequency-polynomial} and the Gaussian
give total error at most
\[
                 4\,4^{2g}e^{-\pi A_0/4}\le e^{-30D}.
\]
For $b\ne0$, use
\[
 (1+y)^{2c}e^{-\pi y^2/2}
                    \le(4c)^c e^{-\pi y^2/4}
\]
and the uniform affine-lattice bound.  Summing the frequency
shells, whose dimension is $d\le d_*$, gives error at most
\begin{align*}
 2(80qD^{3/2})^D
   \sum_{b\ne0}e^{-\pi q\|b\|^2/2}
                            (\|b\|/\lambda_L)^D
 &\le4(120qD^{3/2})^D4^d e^{-\pi B_0/2}\\
 &\le e^{-30D}.
\end{align*}
For the first inequality use the shells $j\ge2$ and
\[
 (j+2)^d((j+1)/2)^D
 \le4^d(3/2)^D e^{(d+D)(j^2-4)/8}.
\]
The resulting geometric sum is less than two since
$B_0\ge d+D+1$.  For the second inequality substitute
\eqref{eq:thin-jet-parameters-b}; its logarithmic margin is at
least $30D$, using $\log120<5$, $\log4<1.4$ and $\pi>3$.

Accordingly, every selected column has total nonidentity error,
in the indicated tensor norm, at most
\begin{equation}\label{eq:thin-column-error}
                            \delta=2e^{-30D}\le(8D)^{-1}.
\end{equation}
The same upper bound applies to the unused zero-degree column
by omitting its polynomial factors.  Polarizing the product of
$c$ determinants expresses its contraction with symmetric
column tensors as an average of products of $c$ determinants.
Each factor is bounded by $\Delta$.  Thus the relative error in
an $i$-column block is at most $(1+\delta)^i-1<1/2$, since
$i\le D$.  Its jet is at least $\Delta^c/2$ in absolute value.
Each unused column has value at least $1/2$.

\emph{The original $L^2$ norm.}
The periodization also has a uniform norm bound; a pointwise
estimate alone would not suffice.  Absolute convergence of the
unfolding can be checked before using orthogonality.  Write
$N(h)=\int_{V/I}|h|^2e^{-\pi q\|z\|^2}$ and let $T_\gamma$
denote the translation in \eqref{eq:thin-periodization}.
For a toric column, completing the square gives exactly
\[
 \frac{\int_{V/I}|T_\gamma h\,\overline h|
                            e^{-\pi q\|z\|^2}}{N(h)}
       =e^{-\pi q\|b\|^2/2-\pi q\|a\|^2/4}.
\]
For an abelian column the same ratio is at most this Gaussian
times $\sum_{l=0}^c(\pi q|H(a,\xi)|^2/4)^l/l!$.
Indeed translate the $A$ coordinate by $a/2$, bound the product
of its two polynomial factors by
$(|H(z,\xi)|^2+|H(a,\xi)|^2/4)^c$, and integrate its monomials.
These Gaussian-polynomial sums converge by packing, uniformly
on the finite-volume real $U$ domain.  Thus the absolute double
sum can be unfolded to $V/I$.

Now integrate first over $U/I$.  The characters are
$u-qb$, so orthogonality kills every term with $b\ne0$.
The remaining terms are ordinary Gaussian overlaps in $A$.
For a toric column their normalized absolute value is
$e^{-\pi q\|a\|^2/2}$.  For an abelian column it is
\[
 e^{-\pi q\|a\|^2/2}
       \left|\sum_{l=0}^c\binom cl
              \frac{(-\pi q|H(a,\xi)|^2)^l}{l!}\right|
 \le e^{-\pi y^2/2+2\sqrt{\pi c}\,y}.
\]
This formula follows by integrating monomials in the one
coordinate $\xi$; its Gaussian moment is
$c!/(\pi q)^c$.  The preceding zero-frequency tail estimate
shows that the sum of the absolute nonidentity overlaps is
less than one.

The squared norm of the unperiodized toric function on $V/I$ is
\[
 \operatorname{covol}(I)(2q)^{-d/2}q^{-(g-d)}
                         e^{2\pi\|u\|^2/q}.
\]
Its normalization $(2\pi i)^{-c}$ only decreases this.
For the abelian polynomial the additional factor is
$c^c/(\pi^c c!)\le(e/\pi)^c\le1$.
By \eqref{eq:thin-isotropic-covolume} and $f\Omega^2\le R^7$,
\[
 \frac{\Omega\theta}{\sqrt q}\le t_0^{-1/2}\le1.
\]
The prefactor of the Gaussian norm is therefore at most one.
We have proved
\begin{equation}\label{eq:thin-column-l2}
                  \|\Theta_h\|_2^2
                           \le2e^{2\pi\|u\|^2/q},
\end{equation}
where $u=0$ for an abelian or unused column.

Take the tensor product of all the selected columns.  The
determinant blocks and unused columns give
$|P(\partial)F(0)|\ge2^{-n}4^{-D}$, because their number is at
most $n$.  Equation~\eqref{eq:thin-column-l2} gives
$\|F\|_2^2\le2^n e^{2\pi nC\phi/q}$ by
\eqref{eq:thin-frequency-radius}.  This proves
\eqref{eq:thin-jet-lower}.  Finally $K_6\ge2$ and $r\ge2$
give the asserted comparison with the prescribed bound.
\end{proof}

\paragraph{Return to the actual cusp sections.}
The required polynomial really has disjoint column sets.  For
the original antitone exponents, give the minor of size $i$ at
$\tau$ its own $i$ columns and raise it to
$a_{\tau,i-1}-a_{\tau,i}$, omitting zero exponents.  This uses at
most
$\sum_\tau s_\tau(s_\tau+1)/2\le fr(r+1)/2\le n$ columns.
Its row degrees telescope to the prescribed exponents.  Different
embeddings, including conjugate embeddings, use disjoint columns,
so every harmonic contraction is zero.  The leading minors give
the required highest-weight relations, and their product is
nonzero with integral coefficients.  Its degree is $D$, and its
coefficient sum is at most $r^D$.  This is the actual normalized
block polynomial; the construction does not use overlapping
columns of nested minors.

The principally polarized theta section has original probability
norm one.  Pulling it back by the prescribed half-level generator
and translating through the finite half-level group gives an
orthonormal basis of $\mathcal H_M$; on $n$ copies its cardinality
is $M^{gn}$.  Indeed $\alpha\bar\alpha=M$, and the original
semicharacter compatibility makes pullback by $\alpha$ a section
of this very line.  Pullback preserves its probability integral
by Haar invariance under the finite isogeny; the translations are
unitary, so every vector has squared norm exactly one.
The half-level ideal has norm $M^f$, and the original projective
rank-$r$ lattice has quotient cardinality $M^{fr}$.
To check completeness, the translations are distinct
characters for the complementary finite Heisenberg subgroup, so
they are orthogonal; their number is $M^g=\dim\mathcal H_M$.
The dimension identity follows directly by writing a theta section
in Fourier series along a symplectic integral basis: its
coefficients are determined freely by the $M^g$ residue classes.
The original normalized theta dictionary identifies these basis
vectors with its half-level rules.  Before taking absolute values
the evaluation factor is
$M^{-D}\kappa(\sqrt{2\xi_R})^D$, with
$|\kappa|^2=M^D$.  Thus its squared absolute value is exactly
$(2\xi_R/M)^D$.  The Lie-norm factor is
$(2\pi\xi_R)^{-D/2}$.  Hence Parseval and
Theorem~\ref{thm:actual-thin-determinant-jet} give, for the original
rules $s_\zeta$,
\begin{equation}\label{eq:thin-actual-theta-sum}
 M^D\sum_{\zeta}
       \bigl((2\pi\xi_R)^{-D/2}\|s_\zeta(x)\|\bigr)^2
             \ge\pi^{-D}\,2^{-3n}16^{-D}
                           e^{-2\pi nC\phi(x)/M}.
\end{equation}
The integrality, harmonicity, cuspidality and original evaluation
normalization of these rules were established in the actual theta
construction.  The auxiliary section just constructed selects a
nonzero value among them; it does not replace their arithmetic
lattice or their metric.

\begin{theorem}[The original thin-region Bergman bound]
\label{thm:actual-thin-bergman}
Retain the original sized good cover, anchored divided weight
$u_1$, flat-weight comparison and their numerical hypotheses.
At every coefficient embedding and every point of an original
component $Z$ satisfying
$M_{\mathrm{deep}}\lambda(x)^2<b$, one has
\begin{equation}\label{eq:actual-thin-bergman}
                  B_Z(x)\ge
                  \exp\left(-\frac{2f}{625R^2\lambda(x)^2}\right).
\end{equation}
Here $\lambda(x)^2$ is exactly the original first minimum
squared, and $B_Z$ uses the original component probability norm.
\end{theorem}

\begin{proof}
We spell out the numerical and metric interfaces.  On the original
minimal level $M$ the numerical hypotheses give
\[
 b\ge1000,\quad gn\le b,\quad r\le n,\quad R^7\le b,
 \quad b^{16}\le M\le b^{70},\quad 1\le D\le b^2,
 \quad M_{\mathrm{deep}}\ge bM.
\]
They give \eqref{eq:thin-jet-parameters-a}--\eqref{eq:thin-jet-parameters-b}
for this very period torus.  The deep-to-minimal-level dictionary
preserves the first minimum and the point metric.  Thus a point
under consideration has $M\lambda(x)^2<1$ and $\lambda(x)<\theta$, and
\[
                     \phi(x)=\frac{2f}{R^2\lambda(x)^2}
                                      \ge b^{14}.
\]
The normalized integral block polynomial has coefficient sum at
most $e^{Db}$.  The already proved actual raw-theta upper bound,
with $L=3\log b$ and size bound $b^2$, gives
\[
 (2\pi\xi_R)^{-D/2}\sup_z\|s_\zeta(z)\|
   \le\exp\bigl(Db+9000(3fr\log b)^2b^2\bigr).
\]
There is no determinant-index factor for this original principal
theta family.  Averaging \eqref{eq:thin-actual-theta-sum} over
its $M^{gn}$ indices gives at least one nonzero value with
\[
 -\log\bigl((2\pi\xi_R)^{-D/2}\|s_\zeta(x)\|\bigr)
 \le3b^3D+\frac{\phi(x)}{2500}+(D+gn)\log M.
\]
Indeed the nonconstant coefficient is
$\pi nC/M\le1/2500$: use $n,m\le b$, $d_*\le2b$,
$R^7\le b$ and $M\ge b^{16}$.
The remaining logarithms from
\eqref{eq:thin-actual-theta-sum} are less than $3b^3D$.

The total terms not containing $\phi$ are at most $10^5b^6$,
using $fr\le b$, $\log b\le b$ and $\log M\le70b$.
Since $b\ge1000$ and $\phi\ge b^{14}$, they are at most
$\phi/2500$.  The same original cusp section therefore satisfies
\[
                    \|s_\zeta(x)\|
                       \ge e^{-\phi(x)/1250}
                                      \sup_z\|s_\zeta(z)\|.
\]
Pull it to the prescribed cover.  The established flat-weight
comparison from the reference theta weight to $u_1$ gives a
section $g$ supported on the same $Z$ and a constant $K>0$ with
$\|g(z)\|=K\|s_\zeta(z)\|$ for every $z$ on $Z$.
It is the product of the already constructed key and parallel
twisting sections.  Thus it preserves the ratio on this very
component.  Take this section supported on $Z$.
Its original probability $L^2$ norm is at most its supremum norm.
Squaring the ratio proves \eqref{eq:actual-thin-bergman}.
\end{proof}

The theorem supplies the lower thin-region norm used in the
geometric section argument.  The construction of the original
numerical family remains a separate obligation wherever it is
still an input.  Point heights and the final arithmetic section
remainder are not supplied by this analytic theorem.
